\documentclass[11pt]{article}

\usepackage[utf8]{inputenc}
\usepackage[T1]{fontenc}
\usepackage{lmodern}
\usepackage[margin=1in]{geometry}
\usepackage{microtype}
\usepackage{amsmath,amssymb,amsthm,mathtools,mathrsfs,bbm}
\usepackage{enumitem}
\usepackage{xcolor}
\usepackage{hyperref}
\usepackage[nameinlink,noabbrev]{cleveref}
\usepackage{tikz-cd}
\usepackage{stmaryrd}
\usepackage{csquotes}
\usepackage[backend=biber,style=alphabetic,maxbibnames=99]{biblatex}
\addbibresource{modular_homotopy_holography_refs.bib}

\hypersetup{
  colorlinks=true,
  linkcolor=blue!50!black,
  citecolor=blue!50!black,
  urlcolor=blue!50!black,
  pdftitle={A Modular Homotopy Package for Beta-Gamma, Kodaira--Spencer, and W3 Holography},
  pdfauthor={OpenAI}
}

\setlist[itemize]{leftmargin=2.1em}
\setlist[enumerate]{leftmargin=2.4em}

\newtheorem{theorem}{Theorem}[section]
\newtheorem{proposition}[theorem]{Proposition}
\newtheorem{corollary}[theorem]{Corollary}
\newtheorem{lemma}[theorem]{Lemma}
\theoremstyle{definition}
\newtheorem{definition}[theorem]{Definition}
\newtheorem{conjecture}[theorem]{Conjecture}
\newtheorem{example}[theorem]{Example}
\newtheorem{question}[theorem]{Question}
\theoremstyle{remark}
\newtheorem{remark}[theorem]{Remark}

\newcommand{\A}{\mathcal A}
\newcommand{\B}{\mathcal B}
\newcommand{\C}{\mathcal C}
\newcommand{\D}{\mathcal D}
\newcommand{\E}{\mathbb E}
\newcommand{\F}{\mathcal F}
\newcommand{\G}{\mathbb G}
\newcommand{\Hh}{\mathcal H}
\newcommand{\Kk}{\mathfrak K}
\newcommand{\M}{\mathcal M}
\newcommand{\ModEnv}{\mathfrak M}
\newcommand{\MC}{\operatorname{MC}}
\newcommand{\End}{\operatorname{End}}
\newcommand{\Hom}{\operatorname{Hom}}
\newcommand{\Sym}{\operatorname{Sym}}
\newcommand{\Conf}{\operatorname{Conf}}
\newcommand{\Ob}{\operatorname{Ob}}
\newcommand{\Tr}{\operatorname{Tr}}
\newcommand{\Vir}{\operatorname{Vir}}
\newcommand{\Ran}{\operatorname{Ran}}
\newcommand{\cl}{\mathrm{cl}}
\newcommand{\hol}{\mathrm{hol}}
\newcommand{\modu}{\mathrm{mod}}
\newcommand{\eff}{\mathrm{eff}}
\newcommand{\defect}{\mathrm{def}}
\newcommand{\op}{\mathrm{op}}
\newcommand{\gr}{\mathrm{gr}}
\newcommand{\Fact}{\mathrm{Fact}}
\newcommand{\Bun}{\mathrm{Bun}}
\newcommand{\Spec}{\mathrm{Spec}}
\newcommand{\id}{\mathrm{id}}
\newcommand{\one}{\mathbbm{1}}
\newcommand{\bbeta}{\beta}
\newcommand{\bgamma}{\gamma}

\title{\vspace{-1em}\textbf{A Modular Homotopy Package for $\bbeta\bgamma$, Kodaira--Spencer, and $W_3$ Holography}\\[0.45em]
\large A theorem--definition--conjecture note with new modular objects}
\author{}
\date{\today}

\begin{document}
\maketitle

\begin{abstract}
We isolate a genus-zero algebraic package that is already present in boundary chiral algebras, twisted holography, higher perturbative operations, and line-operator OPE, and we formulate a precise modular completion problem for it. The established part is assembled from the bulk/boundary formalism of Costello--Dimofte--Gaiotto, the holographic KK-reduction package of Zeng, the higher-operation formalism of Gaiotto--Kulp--Wu, the Poisson-vertex-to-3d-gauge correspondence of Khan--Zeng, and the line-operator/dg-shifted-Yangian framework of Dimofte--Niu--Py. The new part of the note is a collection of explicit modular objects: a modular dg-shifted Yangian, a holographic tautological action, logarithmic clutching kernels built from logarithmic Fulton--MacPherson geometry, a modular descendant potential, and a modular defect index. The three worked examples are the free $\bbeta\bgamma$ system, Kodaira--Spencer holography, and the $W_3$ Poisson vertex algebra. Throughout, statements marked \emph{Theorem} or \emph{Proposition} are either source-backed or formal syntheses of source-backed results, whereas statements marked \emph{Definition} or \emph{Conjecture} are proposed new constructions.
\end{abstract}

\tableofcontents

\section{Scope and notation}

Let $X$ be a smooth complex curve and let $T$ be a perturbative holomorphic-topological quantum field theory on $\mathbb R_t \times X$. A boundary condition is denoted by $B$. We write
\[
A_T
\]
for the bulk $A_\infty$-chiral algebra, $A_T^!$ for a completed Koszul dual algebra, and
\[
\C_T := A_T^!\text{-mod}
\]
for the category of perturbative line operators whenever the line-operator formalism is available. Stable curves are written $\Sigma_{g,n}$, and $\overline{\M}_{g,n}$ denotes the Deligne--Mumford stack of stable $n$-pointed genus-$g$ curves.

The aim is to turn the familiar genus-zero package
\[
(A_T,\,A_T^!,\,\C_T,\,r_T(z),\,m_2,m_3,\dots)
\]
into a modular object that simultaneously sees stable-curve gluing, higher operations, and all-genera holographic observables.

\begin{remark}[Status convention]
The literature already provides the local genus-zero algebraic pieces. The full modular completion proposed below is not presently a theorem in this level of generality. The note is written so that the boundary between established and conjectural structure is explicit at every step.
\end{remark}

\section{The genus-zero package}

\begin{theorem}[Genus-zero holographic package, assembled from existing results]\label{thm:genus-zero-package}
Let $T$ be a perturbative holomorphic-topological theory of the type considered in \cite{CDG2023,Zeng2024,GKW2025,KZ2025,DNP2025}. Then the following genus-zero data are available.
\begin{enumerate}[label=\textup{(\arabic*)}]
  \item There is a commutative bulk chiral algebra $A_T$ endowed with a cohomological degree $-1$ Poisson bracket. For 3d holomorphic twists of $\mathcal N=2$ theories there is in addition a secondary stress tensor implementing the holomorphic translation operator.\vspace{0.15em}
  \item For each boundary condition $B$, the boundary algebra $A_{\partial,T}(B)$ is a chiral algebra equipped with a bulk--boundary map
  \[
  \beta_B:A_T \longrightarrow Z\bigl(A_{\partial,T}(B)\bigr),
  \]
  so the boundary algebra is a module over the bulk algebra. In general higher $A_\infty$-type structure appears on the boundary side.\vspace{0.15em}
  \item In the KK-reduction setup of twisted and celestial holography, one may choose a boundary condition $B_{\mathrm{univ}}$ for which the boundary chiral algebra coincides with the universal defect chiral algebra of the higher-dimensional theory.\vspace{0.15em}
  \item The deformation, anomaly, and generalized OPE data of perturbative holomorphic-topological theories are controlled by higher operations satisfying quadratic identities, physically interpreted as a Wess--Zumino consistency package.\vspace{0.15em}
  \item Perturbative line operators are modeled by modules for a Koszul-dual $A_\infty$-algebra $A_T^!$, and their OPE is controlled by a Maurer--Cartan kernel
  \[
  r_T(z) \in A_T^! \widehat\otimes A_T^!((z^{-1}))
  \]
  satisfying an $A_\infty$ version of the Yang--Baxter constraint. In the language of \cite{DNP2025}, this package is part of a dg-shifted Yangian.
\end{enumerate}
\end{theorem}

\begin{proof}
Item \textup{(1)} and the bulk--boundary formalism in \textup{(2)} are established in \cite{CDG2023}. Item \textup{(3)} is the boundary-chiral-algebra/universal-defect identification of \cite{Zeng2024}. Item \textup{(4)} is the higher-operation package of \cite{GKW2025}. Item \textup{(5)} is the line-operator and dg-shifted-Yangian formalism developed in \cite{DNP2025}. The theorem is a formal synthesis of these source-backed statements.
\end{proof}

\begin{definition}[Genus-zero package]\label{def:genus-zero-package}
For a theory $T$ with boundary condition $B$, define
\[
\G_0(T;B)
:=
\bigl(A_T,\,\{-,-\}_T,\,A_{\partial,T}(B),\,\beta_B,\,A_T^!,\,\C_T,\,r_T(z),\,\{m_n^T\}_{n\ge 2}\bigr),
\]
where $\{m_n^T\}_{n\ge 2}$ denotes the genus-zero higher operations furnished by perturbation theory or homotopy transfer.
\end{definition}

\section{The modular completion problem}

\begin{definition}[Full modular homotopy type]\label{def:full-modular-homotopy-type}
A \emph{full modular homotopy type} for $T$ based at $B$ is a completion of $\G_0(T;B)$ by multilinear maps
\[
\mu^T_{g,n}:
C_\bullet\bigl(\overline{\M}_{g,n+1}\bigr)
\otimes
(A_T^!)^{\widehat\otimes n}
\longrightarrow
A_T^!,
\qquad 2g-2+n>0,
\]
subject to the following conditions.
\begin{enumerate}[label=\textup{(\alph*)}]
  \item \textbf{Genus-zero recovery:} the operations $\mu_{0,n}^T$ recover the local higher operations $m_n^T$ after restriction to the genus-zero boundary strata.
  \item \textbf{Stable-graph gluing:} the collection $\{\mu^T_{g,n}\}$ is compatible with the stable-graph composition law on the modular operad $C_\bullet(\overline{\M}_{\bullet,\bullet})$.
  \item \textbf{Line-kernel refinement:} the genus-zero Maurer--Cartan element extends to a formal expansion
  \[
  R_T^{\modu}(z;\hbar)
  :=
  \sum_{g\ge 0}\hbar^{2g}r_{T,g}(z),
  \qquad
  r_{T,0}(z)=r_T(z),
  \]
  with
  \[
  r_{T,g}(z)\in H^\bullet\bigl(\overline{\M}_{g,2}\bigr)\otimes A_T^!\widehat\otimes A_T^!((z^{-1})).
  \]
\end{enumerate}
\end{definition}

\begin{conjecture}[Modular envelope conjecture]\label{conj:modular-envelope}
For the classes of theories appearing in \cref{thm:genus-zero-package}, every genus-zero package $\G_0(T;B)$ admits a full modular homotopy type $\ModEnv(T;B)$.
\end{conjecture}

\begin{remark}
The conjecture is deliberately stronger than the mere existence of genus-by-genus correlation functions. It asks for an actual modular operadic object whose gluing law controls higher operations, line OPE, and defect observables in a single package.
\end{remark}

\begin{definition}[Modular state spaces and partition function]\label{def:state-spaces}
Assume $\ModEnv(T;B)$ exists. For line objects $L_1,\dots,L_n\in\C_T$, define the \emph{modular state space}
\[
\Hh_{g,n}^{T;B}(L_1,\dots,L_n)
:=
\int_{\Sigma_{g,n}} \Fact_{T;B}[L_1,\dots,L_n],
\]
namely the factorization or chiral homology of the local operator package over the stable pointed curve. The associated genus-graded partition function is
\[
\mathcal Z_{T;B}(\hbar,\mathbf t)
:=
\sum_{g,n\ge 0}
\hbar^{2g-2+n}\frac{1}{n!}
\chi_{\gr}\bigl(\Hh_{g,n}^{T;B}(\mathbf t)\bigr).
\]
\end{definition}

\begin{remark}
The factorization-homology viewpoint is natural because it is literally the operation of integrating local algebraic data over a manifold, and because it comes with local-to-global, excision, and Poincar\'e/Koszul duality features \cite{AF2015,AF2019,AFPrimer,BD2004,BrochierWoike2022}. In our context, the manifold is replaced by a stable complex curve together with marked points and defects.
\end{remark}

\section{New modular constructions and new quantum quantities}\label{sec:new-objects}

This section is the genuinely forward-driving part of the note. The objects below are proposed here as precise modular enhancements of the known genus-zero package.

\subsection{The modular dg-shifted Yangian}

\begin{definition}[Modular dg-shifted Yangian]\label{def:modular-yangian}
Let $\ModEnv(T;B)$ be a full modular homotopy type. The \emph{modular dg-shifted Yangian} of $(T,B)$ is the stable-graph completion
\[
\mathbb Y_T^{\modu}
:=
\Bigl(A_T^!,\,\{\mu^T_{g,n}\}_{2g-2+n>0},\,\{r_{T,g}(z)\}_{g\ge 0}\Bigr)
\]
viewed as a completed colored modular algebraic structure whose unary and binary operations are generated by the stable-curve classes and the genus-refined kernels $r_{T,g}(z)$.
\end{definition}

\begin{conjecture}[Modular Yang--Baxter package]\label{conj:modular-yang-baxter}
The kernels $r_{T,g}(z)$ satisfy a stable-graph refinement of the $A_\infty$ Yang--Baxter constraint of \cite{DNP2025}. Equivalently, the total kernel $R_T^{\modu}(z;\hbar)$ is a Maurer--Cartan element in a pronilpotent dg Lie algebra built from the boundary stratification of $\overline{\M}_{g,n}$.
\end{conjecture}

\begin{remark}[Why this is new]
The usual dg-shifted Yangian is genus-zero and local. The object $\mathbb Y_T^{\modu}$ is intended to remember both line-operator OPE and stable-curve gluing in one envelope. In practice, it should be the correct home for all-genera defect OPE coefficients.
\end{remark}

\subsection{A holographic tautological action}

\begin{definition}[Universal correlator class]\label{def:universal-correlator-class}
Suppose $\ModEnv(T;B)$ exists. For each $(g,n)$, define the \emph{universal correlator class}
\[
\mathbb U^{T;B}_{g,n}
\in
H^\bullet\bigl(\overline{\M}_{g,n};\End(\Hh_{g,n}^{T;B})\bigr)
\]
by packaging the descendants produced by the operations $\mu^T_{g,n}$ and the factorization-homology pushforward over the universal $n$-pointed genus-$g$ curve.
\end{definition}

\begin{definition}[Holographic tautological action]\label{def:tautological-action}
The \emph{holographic tautological action} is the map
\[
\Theta^{T;B}_{g,n}:
R^\bullet\bigl(\overline{\M}_{g,n}\bigr)
\longrightarrow
\End\bigl(\Hh_{g,n}^{T;B}\bigr),
\qquad
\alpha \longmapsto \int_{\overline{\M}_{g,n}} \alpha \smile \mathbb U_{g,n}^{T;B},
\]
where $R^\bullet(\overline{\M}_{g,n})$ denotes the tautological ring.
\end{definition}

\begin{remark}
This is an explicit new object. It upgrades numerical descendant integrals to an operator-valued action of the tautological ring on modular state spaces. If it exists, one gets a direct bridge from tautological relations on $\overline{\M}_{g,n}$ to protected operator identities in holography.
\end{remark}

\subsection{Logarithmic clutching kernels}

\begin{definition}[Logarithmic clutching kernel]\label{def:log-clutching-kernel}
Fix a pair $(X,D)$ with $X$ smooth and $D\subset X$ a normal-crossings divisor. Let
\[
FM_n(X\mid D)
\]
be the logarithmic Fulton--MacPherson compactification of $n$ ordered points on $X\setminus D$ in the sense of \cite{Mok2025}. For stable data $(g_1,n_1+1)$ and $(g_2,n_2+1)$, define a \emph{logarithmic clutching kernel} to be a proper correspondence class
\[
\Kk^{\log}_{g_1,n_1;g_2,n_2}
\in
\mathrm{Corr}\Bigl(
FM_{n_1+1}(X\mid D)\times FM_{n_2+1}(X\mid D),
FM_{n_1+n_2}(X\mid D)
\Bigr)
\]
that implements the gluing of the last marked point on the first factor with the last marked point on the second factor.
\end{definition}

\begin{conjecture}[Logarithmic clutching from degeneration geometry]\label{conj:log-clutching}
The degeneration formula of \cite{Mok2025} canonically produces logarithmic clutching kernels compatible with the stable-graph boundary stratification of $\overline{\M}_{g,n}$. In particular, every rigid combinatorial type of logarithmic FM degeneration should determine a correspondence component of $\Kk^{\log}_{g_1,n_1;g_2,n_2}$.
\end{conjecture}

\begin{remark}[Conceptual cohesion]
This is the geometric hinge of the program: stable-curve gluing lives on the Deligne--Mumford side, while configuration-space compactification lives on the logarithmic FM side. The kernel $\Kk^{\log}$ is designed to make them interact functorially.
\end{remark}

\subsection{Descendant potentials and defect indices}

\begin{definition}[Modular holographic descendant potential]\label{def:descendant-potential}
Let $\{O_a\}_{a\in I}$ be a homogeneous basis of bulk or defect insertions. The \emph{modular holographic descendant potential} is
\[
\F_{T;B}(\hbar,\mathbf t,\mathbf s)
:=
\sum_{g,n\ge 0}\;
\sum_{a_1,\dots,a_n\in I}\;
\sum_{k_1,\dots,k_n\ge 0}
\frac{\hbar^{2g-2+n}}{n!}
\Bigl\langle \prod_{i=1}^n \tau_{k_i}(O_{a_i})\Bigr\rangle^{T;B}_{g,n}
\prod_{i=1}^n t_{a_i}s_{k_i},
\]
where the descendants are defined through the tautological action of \cref{def:tautological-action}.
\end{definition}

\begin{definition}[Modular defect index]\label{def:defect-index}
The \emph{modular defect index} is the graded trace
\[
\mathfrak I^{\defect}_{T;B}(q,\hbar)
:=
\sum_{g,n\ge 0}
\hbar^{2g-2+n}
\Tr_{\Hh_{g,n}^{T;B}}
\bigl((-1)^F q^{L_0}\bigr),
\]
when the $L_0$-grading is defined on the relevant state spaces.
\end{definition}

\begin{conjecture}[Tautological recursion for holographic descendants]\label{conj:tautological-recursion}
The potential $\F_{T;B}$ satisfies all tautological relations imposed by the image of $R^\bullet(\overline{\M}_{g,n})$ under $\Theta^{T;B}_{g,n}$. In favorable cases this should produce topological-recursion-type identities for modular holographic amplitudes.
\end{conjecture}

\section{Worked example I: the beta-gamma system}\label{sec:beta-gamma}

\begin{proposition}[Interval compactification produces the full $\bbeta\bgamma$ algebra]\label{prop:beta-gamma-interval}
Consider the holomorphic twist of a 3d $\mathcal N=2$ free chiral multiplet on $[0,1]\times \mathbb C$ with Neumann boundary conditions at both ends. Let $\phi$ denote the boundary boson and let $\psi^{(1)}$ be the first descendant of the fermionic bulk field. Then the interval theory contains fields
\[
\bbeta(z):=\phi(z),
\qquad
\bgamma(z):=\int_0^1 \psi^{(1)},
\]
with operator product
\[
\bbeta(z)\bgamma(0)\sim \frac{1}{z}.
\]
Hence the interval compactification realizes the full $\bbeta\bgamma$ chiral algebra rather than a single boundary half of it.
\end{proposition}

\begin{proof}
This is the interval-compactification analysis of \cite[\S4.2]{CDG2023}. The field $\bgamma$ is $Q$-closed by Stokes' theorem because the bulk fermion vanishes on both Neumann boundaries. The pole in the OPE is induced by the bulk Poisson bracket between $\phi$ and $\psi$.
\end{proof}

\begin{definition}[The free $\bbeta\bgamma$ modular test package]\label{def:beta-gamma-modular-package}
Let
\[
V_{\bbeta\bgamma}
=
\Sym\bigl(\mathbb C[\partial]\bbeta\oplus \mathbb C[\partial]\bgamma\bigr),
\qquad
\{\bbeta_\lambda \bgamma\}=1.
\]
The corresponding free modular test object is the tuple
\[
\ModEnv_{\bbeta\bgamma}^{\mathrm{free}}
:=
\bigl(V_{\bbeta\bgamma},\,V_{\bbeta\bgamma}^!,\,\Hh_{g,n}^{\bbeta\bgamma},\,\Theta^{\bbeta\bgamma}_{g,n},\,\F_{\bbeta\bgamma},\,\mathfrak I_{\bbeta\bgamma}^{\defect}\bigr),
\]
with higher local operations taken to be strict and all higher-genus complexity pushed into factorization homology and the tautological action.
\end{definition}

\begin{definition}[The $\bbeta\bgamma$ determinant line]\label{def:beta-gamma-det-line}
Fix a conformal weight parameter $\lambda$ and an auxiliary line bundle $L$ on $\Sigma_g$. Define the \emph{$\bbeta\bgamma$ determinant line}
\[
\mathfrak D^{\bbeta\bgamma}_{g,\lambda}(L)
:=
\det R\Gamma\bigl(\Sigma_g,K_{\Sigma_g}^{\lambda}\otimes L\bigr)^{-1}
\otimes
\det R\Gamma\bigl(\Sigma_g,K_{\Sigma_g}^{1-\lambda}\otimes L^{\vee}\bigr)^{-1}.
\]
\end{definition}

\begin{conjecture}[Free-field realization of modular state spaces]\label{conj:beta-gamma-det-line}
For the free $\bbeta\bgamma$ theory, the state space $\Hh_g^{\bbeta\bgamma}$ is governed by the determinant geometry of \cref{def:beta-gamma-det-line}. More precisely, the genus-$g$ partition function should be expressible as a canonical section of the determinant line obtained by integrating the Dolbeault complexes underlying $\bbeta$ and $\bgamma$.
\end{conjecture}

\begin{remark}
This conjecture packages a familiar free-field heuristic into a modular-homotopy statement: the local OPE is free, but the global all-genera content is encoded by determinant lines on curves. In this sense the free $\bbeta\bgamma$ theory is the cleanest testing ground for the modular envelope conjecture.
\end{remark}

\section{Worked example II: Kodaira--Spencer holography}\label{sec:KS}

\begin{example}[Large-$N$ boundary algebra and universal defect generators]\label{ex:KS-boundary-algebra}
In the twisted-holography package discussed in \cite{Zeng2024,KZ2025}, the large-$N$ open-side chiral algebra is built from a BRST system with fields
\[
(b,c)\ \text{in}\ \mathfrak{gl}_N,
\qquad
(Z_1,Z_2)\ \text{an adjoint $\bbeta\bgamma$ system},
\qquad
(I,J)\ \text{a fundamental $\bbeta\bgamma$ system},
\]
and BRST operator
\[
Q_{\mathrm{BRST}}
=
\oint \Bigl(
\mathrm{Tr}(cZ_1Z_2)
+
\mathrm{Tr}(IcJ)
+
\frac{1}{2}\mathrm{Tr}(bc^2)
\Bigr).
\]
Its large-$N$ BRST cohomology contains towers of generators
\[
A^{(n)},\ B^{(n)},\ C^{(n)},\ D^{(n)},\ E_t^{(n)},
\]
where the first four towers match the Kodaira--Spencer gravity sector and the $E$-tower matches the holomorphic Chern--Simons sector. This boundary algebra is the natural candidate for the universal defect algebra of the Kodaira--Spencer side.
\end{example}

\begin{definition}[Kodaira--Spencer universal defect algebra]\label{def:KS-universal-defect}
Define
\[
U_{\mathrm{KS}}
:=
A_{\partial,\mathrm{KS}}(B_{\mathrm{univ}}),
\qquad
U_{\mathrm{KS}}^{\cl}
:=
\langle A^{(n)},B^{(n)},C^{(n)},D^{(n)}\rangle,
\qquad
U_{\mathrm{KS}}^{\mathrm{gauge}}
:=
\langle E_t^{(n)}\rangle.
\]
The corresponding modular state spaces are
\[
\Hh_{g,n}^{\mathrm{KS}}:=\int_{\Sigma_{g,n}}U_{\mathrm{KS}},
\qquad
\Hh_{g,n}^{\cl,\mathrm{KS}}:=\int_{\Sigma_{g,n}}U_{\mathrm{KS}}^{\cl}.
\]
\end{definition}

\begin{conjecture}[Open--closed modular splitting for Kodaira--Spencer holography]\label{conj:KS-splitting}
There is a functorial completed factorization
\[
\Hh_{g,n}^{\mathrm{KS}}
\simeq
\Hh_{g,n}^{\cl,\mathrm{KS}}
\widehat\otimes
\Hh_{g,n}^{\mathrm{gauge},\mathrm{KS}},
\]
compatible with the holographic bulk--boundary map and with the stable-curve gluing operations $\mu^{\mathrm{KS}}_{g,n}$. Equivalently, the Kodaira--Spencer modular envelope should carry an intrinsic open--closed splitting into a gravitational piece and a gauge-theoretic defect piece.
\end{conjecture}

\begin{remark}
This proposed splitting is one of the central new insights of the note. It gives a mathematically sharp candidate for separating the pure Kodaira--Spencer gravity sector from the holomorphic Chern--Simons sector at all genera while remaining inside a single universal defect algebra.
\end{remark}

\section{Worked example III: the W3 Poisson vertex algebra}\label{sec:W3}

\begin{example}[Classical $W_3$ brackets and the associated 3d action]\label{ex:W3-classical}
Let
\[
\mathcal P_{W_3}
=
\Sym\bigl(\mathbb C[\partial]T\oplus \mathbb C[\partial]W\bigr)
\]
be the classical $W_3$ Poisson vertex algebra. The basic $\lambda$-brackets are
\begin{align*}
\{T_\lambda T\}
&=
\partial T + 2T\lambda + \frac{c}{12}\lambda^3,
\\
\{T_\lambda W\}
&=
\partial W + 3W\lambda,
\\
\{W_\lambda W\}
&=
\frac{c}{360}\lambda^5
+\frac{1}{3}T\lambda^3
+\frac{1}{2}(\partial T)\lambda^2
+\Bigl(\frac{3}{10}\partial^2T+\frac{32}{5c}T^2\Bigr)\lambda
+\frac{1}{15}\partial^3T
+\frac{32}{5c}T\partial T.
\end{align*}
The associated 3d mixed holomorphic-topological gauge theory has fields $(T,\mu)$ and $(W,\chi)$ with interaction terms explicitly determined by the above brackets; see \cite[\S3.4]{KZ2025}.
\end{example}

\begin{proposition}[Topological enhancement in the $W_3$ case]\label{prop:w3-topological-enhancement}
Because the $W_3$ Poisson vertex algebra contains the Virasoro field $T$, the associated 3d holomorphic-topological Poisson sigma model enjoys full topological symmetry enhancement in the sense of \cite{KZ2025}.
\end{proposition}

\begin{proof}
This is a special case of the general criterion in \cite{KZ2025}: the presence of a Virasoro element upgrades the mixed holomorphic-topological symmetry to full topological symmetry.
\end{proof}

\begin{example}[Slab reduction and the central-charge shift]\label{ex:w3-central-charge}
Under slab reduction, the $W_3$ theory becomes a pair of bosonic $\bbeta\bgamma$ systems, one of spins $(2,-1)$ and one of spins $(3,-2)$. Using the standard formula
\[
 c(s)=3(2s-1)^2-1 = 2(6s^2-6s+1),
\]
we obtain
\[
 c(2)=26,
 \qquad
 c(3)=74,
 \qquad
 c_{\mathrm{slab}}=100.
\]
Khan and Zeng therefore predict that the deformation quantization of the classical $W_3$ Poisson algebra should yield the quantum $W_3$ algebra with shifted central charge
\[
 c_{\eff}=c+50.
\]
\end{example}

\begin{definition}[Modular $W_3$ envelope]\label{def:modular-w3-envelope}
The \emph{modular $W_3$ envelope} is the putative tuple
\[
\ModEnv(W_3)
:=
\bigl(W_3^{\mathrm q}(c_{\eff}),\,(W_3^{\mathrm q}(c_{\eff}))^!,\,\mathbb Y_{W_3}^{\modu},\,\Theta_{g,n}^{W_3},\,\F_{W_3},\,\mathfrak I_{W_3}^{\defect}\bigr),
\]
where $W_3^{\mathrm q}(c_{\eff})$ denotes the quantum $W_3$ vertex algebra at the shifted central charge $c_{\eff}=c+50$.
\end{definition}

\begin{conjecture}[Higher-spin modular control]\label{conj:w3-higher-spin}
The modular dg-shifted Yangian $\mathbb Y_{W_3}^{\modu}$ controls the all-genera OPE of $W_3$ line operators, and the tautological action $\Theta_{g,n}^{W_3}$ records a genuine higher-spin response of the modular state spaces to tautological classes on $\overline{\M}_{g,n}$.
\end{conjecture}

\begin{remark}
The $W_3$ case is the first nonlinear test of the program. The $\bbeta\bgamma$ example is free, and the Kodaira--Spencer example is holographically rich but still organized by large-$N$ BRST geometry. The $W_3$ example is where nonlinear higher-spin OPE, stable-curve gluing, and modular line-operator structure genuinely meet.
\end{remark}

\section{A cohesive research program}\label{sec:program}

The constructions above suggest a compact slogan:
\[
\begin{aligned}
\textit{stable-graph summation} &+ \textit{logarithmic configuration geometry}\\[0.2em]
&+ \textit{dg-shifted Yangian line algebra}\\[0.2em]
&= \textit{modular holography}.
\end{aligned}
\]

To make this program mathematically effective, one should try to prove the following three statements.

\begin{question}[Existence]
Can one construct $\ModEnv(T;B)$ directly from the genus-zero package by a bar--cobar or deformation-theoretic completion over the modular operad $C_\bullet(\overline{\M}_{g,n})$?
\end{question}

\begin{question}[Geometry]
Can one realize the gluing operations $\mu_{g,n}^T$ by explicit integrals over logarithmic Fulton--MacPherson spaces and their degeneration correspondences?
\end{question}

\begin{question}[Computation]
Can one compute the first genuinely modular coefficients of $R_T^{\modu}(z;\hbar)$ in the three basic examples: free $\bbeta\bgamma$, Kodaira--Spencer holography, and the $W_3$ Poisson sigma model?
\end{question}

\begin{remark}[What is genuinely novel here]
The note introduces five concrete objects that were not part of the source literature as stated: the modular dg-shifted Yangian $\mathbb Y_T^{\modu}$, the holographic tautological action $\Theta_{g,n}^{T;B}$, the logarithmic clutching kernels $\Kk^{\log}$, the modular descendant potential $\F_{T;B}$, and the modular defect index $\mathfrak I_{T;B}^{\defect}$. Their value is that they turn the vague phrase \emph{all-genera holographic data} into an actual algebraic-geometric package with testable outputs.
\end{remark}

\section{Conclusion}

The established papers already explain how to pass from bulk fields to boundary chiral algebras, from KK reduction to universal defect algebras, from perturbation theory to higher operations, from line operators to a dg-shifted Yangian, and from local algebra to factorization homology. The modular step proposed here is to \emph{bind all of these ingredients to the stable geometry of curves}. In that form, the free $\bbeta\bgamma$ theory becomes a determinant-line test case, Kodaira--Spencer holography becomes an open--closed modular factorization problem, and $W_3$ becomes the first nonlinear higher-spin arena for modular dg-shifted Yangians.

\begingroup\footnotesize\setlength{\bibitemsep}{0pt}\printbibliography\endgroup

\end{document}
